\chapter*{Executive Summary}
\addcontentsline{toc}{chapter}{Executive Summary}

This report summarises my internship at \textbf{FlowGenX AI}, a San
Francisco-based software company building an \emph{Agentic Integration
Platform} for designing, orchestrating, and deploying AI agents and
workflows through a single no-/low-code interface. The internship was
completed as the partial fulfilment of the SWE--420 Internship course
of the Bachelor of Science in Software Engineering programme at IICT,
Shahjalal University of Science and Technology, Sylhet.

I worked at FlowGenX AI as a \emph{Full Stack Software Engineer ---
Generative AI (Intern)} from \textbf{10 November 2025} to
\textbf{02 May 2026}, full-time and remotely from Bangladesh. From
the first week I was placed on the connectivity team and trusted with
production-grade work that ships to real users on the platform.

The bulk of my contribution went to the \texttt{runtime\_connectivity}
project --- FlowGenX's unified Python framework that exposes more than
two hundred external services through a single standardised interface.
Over five and a half months I authored \textbf{178 commits} and
delivered \textbf{more than ninety production-ready connectors}, each
implemented end-to-end against the project's three-tier abstraction,
validated by dedicated integration test scripts, and accompanied by
auto-generated metadata. Beyond feature work I shipped reliability
fixes (OAuth 2.0 refresh, retry logic, encoding fixes) and authored
UI Endpoint Test Reference documentation for many high-value
connectors.

The internship also exposed me to the engineering culture of a
globally distributed startup --- a disciplined feature-branch Git
workflow, peer-reviewed pull requests, automated code-quality
enforcement, and asynchronous collaboration across the United States
and Bangladesh time zones --- and gave me practical exposure to
FastAPI, Apache Airflow, the Model Context Protocol (MCP), and the
Agent-to-Agent (A2A) patterns at the heart of FlowGenX's agentic
stack. The remaining chapters present the company profile, the
engineering process, the internship activities, the skills acquired,
remote-work reflections, and a concluding reflection.
